\DocumentMetadata
  {
    lang=en-US,
    pdfversion=2.0,
    pdfstandard=ua-2,
    tagging=on
  }
\documentclass{article}
\usepackage{attrib}

\title{attrib tagging test}

\begin{document}

\begin{quotation}
[My plays] deal with distress. Some people object to this in my writing.
At a party an English intellectual---so called---asked me why I write always
about distress. As if it were perverse to do so! He wanted to know if my
father had beaten me or my mother had run away from home to give me an
unhappy childhood. I told him no, that I had had a very happy childhood.
Then he thought me more perverse than ever. I left the party as soon as
possible and got into a taxi. On the glass partition between me and the
driver were three signs: one asked for help for the blind, another help for
orphans, and the third for relief for the war refugees. One does not have to
look for distress. It is screaming at you even in the taxis of London.
\attrib{Driver 1961, 24}
\end{quotation}

\begin{verse}
Jacke boy, ho boy newes,\\
\quad the cat is in the well,\\
let us ring now for her Knell,\\
\quad ding dong ding dong Bell.
\attrib{Opie and Opie 1952, 149\footnote{%
Perhaps more familiar is the nursery rhyme that begins ‘‘Ding,
dong, bell, / Pussy’s in the well,’’ which Opie and Opie also cite. It does
not seem clear in that rhyme, however, without contextual knowledge, that the
cat has died and the bell rings its knell.}}
\end{verse}

\end{document}